\chapter*{参考文献}
\addcontentsline{toc}{chapter}{参考文献}
\markboth{参考文献}{参考文献}
\thispagestyle{plain}
\begingroup\small\raggedright\setlength{\parindent}{0pt}
\begin{list}{}{\leftmargin=2.1em \labelwidth=0pt \labelsep=0pt \itemindent=-2.1em \parsep=0pt \itemsep=1.6pt \topsep=4pt \rightmargin=0pt}
\item Aalen OO, Stensrud MJ, Didelez V, Daniel R, Røysland K, Strohmaier S (2020). Time-dependent mediators in survival analysis: Modeling direct and indirect effects with the additive hazards model. Biometrical Journal 62(3):532–549
\item Abadie A (2005). Semiparametric difference-in-differences estimators. The Review of Economic Studies 72(1): 1-19.
\item Abadie A, Chingos MM, West MR (2018). Endogenous stratification in randomized experiments. The Review of Economics and Statistics 100(4): 567–580.
\item Abadie A, Imbens GW (2006). Large sample properties of matched estimates for average treatment effects. Econometrica 74:235-267.
\item Amrhein V, Greenland S, McShane B (2019). Scientists rise up against statistical significance. Nature 567:305-307.
\item Angrist JD, Imbens GW (1995). Two-stage least squares estimation of average causal effects in models with variable treatment intensity. Journal of the American Statistical Association 90:431-442.
\item Angrist JD, Imbens GW, Rubin DB (1996). Identification of causal effects using instrumental variables. Journal of the American Statistical Association 91:444-455.
\item Angrist JD, Krueger AB (1999). Empirical strategies in labor economics. In: Ashenfelter O, Card D, eds. Handbook of Labor Economics 3A, 1277-1366. Elsevier.
\item Angrist JD, Pischke J-S (2009). Mostly Harmless Econometrics: An Empiricist's Companion. Princeton University Press.
\item Avagyan V, Vansteelandt S (2021). Stable inverse probability weighting estimation for longitudinal studies. Scandinavian Journal of Statistics 48(3), 1046– 1067.
\item Baiocchi M, Small D, Lorch S, Rosenbaum P (2010). Building a sttronger instrument in an observational study of perinatal care for premature infants. Journal of the American Statistical Association 105(492): 1285-1296.
\item Baiocchi M, Cheng J, Small D (2014). Instrumental variable methods for causal inference. Statistics in Medicine 33: 2297-2340.
\item Baker SG, Lindeman KS (1994). The paired availability design, a proposal for evaluating epidural analgesia during labor. Statistics in Medicine 13:2269-2278.
\item Baker SG, Kramer BS, Lindeman KL (2016). Latent class instrumental variables. A clinical and biostatistical perspective. Statistics in Medicine 35:147-160 (errata in Statistics in Medicine 2019; 38: 901).
\item Balke A, Pearl J (1994). Probabilistic evaluation of counterfactual queries. In Proceedings of the Twelfth National Conference on Artificial Intelligence, Volume I, pp. 230-237.
\item Balke A, Pearl J (1997). Bounds on treatment effects from studies with imperfect compliance. Journal of the American Statistical Association 92(439):1171-1176.
\item Bang H, Robins JM (2005). Doubly robust estimation in missing data and causal inference models. Biometrics 61: 962-972 (errata in Biometrics 2008;64:650).
\item Barnow SN, Cain GG, Goldberger AS (1980). Issues in the analysis of selectivity bias. In: Evaluation Studies (Vol. 5). Stromsdorfer EW, Farkas G, eds. Beverly Hills, CA: Sage Publication.
\item Berkson J (1946). Limitations of the application of fourfold table analysis to hospital data. Biometrics 2: 47-53.
\item Berkson J (1955). The statistical study of association between smoking and lung cancer. Proceedings of the Staff Meetings of the Mayo Clinic 30: 319-348.
\item Blot WJ, Day NE (1979). Synergism and interaction: are they equivalent? [letter] American Journal of Epidemiology 110:99-100.
\item Blyth CR (1972). On Simpson's paradox and the sure-thing principle. Journal of the American Statistical Association 67:364-66.
\item Bonet B (2001). Instrumentality tests revisited. In: Proceedings of the Seventeenth Conference on Uncertainty in Artificial Intelligence. San Francisco, CA: Morgan Kaufmann, pp. 48-55.
\item Bound J, Jaeger D, Baker R (1995). Problems with instrumental variables estimation when the correlation between the instruments and the endogenous explanatory variables is weak. Journal of the American Statistical Association 90:443-450.
\item Brown LD, Cai T, DasGupta A (2001). Interval estimation for a binomial proportion (with discussion). Statistical Science 16:101-133.
\item Brookhart MA, Schneeweiss S (2007). Preference-based instrumental variable methods for the estimation of treatment effects: assessing validity and interpreting results. International Journal of Biostatistics 3:14.
\item Brookhart MA, Rassen J, Schneeweiss S (2010). Instrumental variable methods in comparative safety and effectiveness research. Pharmacoepidemiology and Drug Safety 19:537-554.
\item Buehler RJ (1982). Some ancillary statistics and their properties. Rejoinder. Journal of the American Statistical Association 77:593-594. Cain LE, Robins JM, Lanoy E, Logan R, Costagliola D, Hernán MA. When to start treatment? A systematic approach to the comparison of dynamic regimes using observational data. International Journal of Biostatistics 2010; 6(2) Article 18.
\item Card D (1990). The Impact of the Mariel Boatlift on the Miami Labor Market. Industrial and Labor Relations Review 43(2):245-257.
\item Card D (1995). Using geographic variation in college proximity to estimate the return to schooling. In: Christofides LN, Grant EK, Swidinsky R, eds. Aspects of labor market behavior: Essays in honour of John Vanderkamp. Toronto, Canada: University of Toronto Press.
\item Casella G, Berger RL (2002). Statistical Inference, 2nd ed. Pacific Grove, CA: Duxbury Press.
\item Cochran WG (1972). Observational studies. In: Bancroft TA, ed. Statistical Papers in Honor of George W Snedecor. Iowa State University Press, pp. 77-90
\item Cole SR, Hernán MA (2002). Fallibility in estimating direct effects. International Journal of Epidemiology 31(1):163-165. Cornfield J, Haenszel W, Hammond EC, Lilienfeld AM, Shimkin MB, Wynder EL (1959). Smoking and lung cancer: Recent evidence and a discussion of some questions. Journal of the National Cancer Institute 22:173-203. Reprinted in International Journal of Epidemiology 2009; 38(5): 1175-1191.
\item Cox DR (1958). Planning of Experiments. New York, NY: John Wiley and Sons.
\item Cox DR, Wermuth N (1999). Likelihood factorizations for mixed discrete and continuous variables. Scandinavian Journal of Statistics 26(2): 209–220.
\item Cui Y, Pu H, Shi X, Miao W, Tchetgen Tchetgen EJ (2024). Semiparametric proximal causal inference. Journal of the American Statistical Association 119(546): 1348-1359.
\item Dahabreh IJ, Hernán (2019). Extending inferences from a randomized trial to a target population. European Journal of Epidemiology 34(8): 719-722. Dahabreh IJ, Petito LC, Robertson SE, Hernán MA, Steingrimsson JA (2020a). Toward causally interpretable meta-analysis: transporting inferences from multiple randomized trials to a new target population. Epidemiology 31(3):334-344.
\item Dahabreh IJ, Robertson SE, Steingrimsson JA, Stuart EA, Hernán MA (2020b). Extending inferences from a randomized trial to a new target population. Statistics in Medicine 39(14):1999-2014.
\item Danaei G, Robins JM, Hu FB, Manson J, Hernán MA (2016). Effect of weight loss on coronary heart disease and mortality in middle-aged or older women: sensitivity analysis for unmeasured confounding by undiagnosed disease. Epidemiology 27(2): 302-310.
\item Davey Smith G, Ebrahim S (2004). Mendelian randomization: prospects, potentials, and limitations. International Journal of Epidemiology 33: 30-42.
\item Dawid AP (1979). Conditional independence in statistical theory (with discussion). Journal of the Royal Statistical Society B 41:1-31.
\item Dawid AP (2000). Causal inference without counterfactuals (with discussion). Journal of American Statistical Association 95: 407-424.
\item Dawid AP (2002). Influence diagrams for causal modelling and inference. International Statistical Review 70: 161-189.
\item Dawid, A. P. (2003). Causal inference using influence diagrams: The problem of partial compliance (with discussion). In: Green PJ, Hjort NL, Richardson S, eds. Highly Structured Stochastic Systems. New York, NY: Oxford University Press, 45-65. de Finetti (1972). Probability, Induction, and Statistics. John Wiley \& Sons.
\item Deaton A (2010). Instruments, randomization, and learning about development. Journal of Economic Literature 48 424-455. Detels R, Muñoz A, McFarlane G, Kingsley LA, Margolick JB, Giorgi J,
\item Schrager L, Phair J, for the Multicenter AIDS Cohorts Study (1998). JAMA 280(17): 1497-1503.
\item Diaz I, Williams N, Hoffman KL, Schenck EJ (2021). Nonparametric causal effects based on longitudinal modified treatment policies. Journal of the American Statistical Association 118(542):846-857.
\item Didelez V (2019). Defining causal mediation with a longitudinal mediator and a survival outcome. Lifetime Data Analysis 25(4):593–610.
\item Didelez V, Sheehan N (2007). Mendelian randomization as an instrumental variable approach to causal inference. Statistical Methods in Medical Research 16: 309-330.
\item Ding P, VanderWeele TJ, Robins JM (2017). Instrumental variables as bias amplifiers with general outcome and confounding. Biometrika 104(2): 291-302.
\item Dorn HF (1953). Philosophy of inferences for retrospective studies. American Journal of Public Health 43: 677-83.
\item Dosemeci M, Wacholder S, Lubin JH (1990). Does nondifferential misclassification of exposure always bias a true effect toward the null value? American Journal of Epidemiology 132:746-748.
\item Earle CC, Tsai JS, Gelber RD, Weinstein MC, Neumann PJ, Weeks JC (2001). Cancer in the elderly: instrumental variable and propensity analysis. Journal of Clinical Oncology 19(4): 1064-1070.
\item Efron B, Hinkley DV (1978). Assessing the accuracy of the maximum likelihood estimator: observed versus expected Fisher information. Biometrika 65: 657-687.
\item Feinstein AR (1971). Clinical biostatistics. XI. Sources of 'chronology bias' in cohort statistics. Clinical Pharmacology and Therapeutics 12(5): 864-79.
\item Fisher RA (1922). On the mathematical foundations of theoretical statistics. Philosophical Transactions of the Royal Society of London. Series A; 222(594–604): 309-368.
\item Flanders WD (2006). On the relation of sufficient component cause models with potential (counterfactual) models. European Journal of Epidemiology 21: 847-853. Flanders WD, Klein M, Darrow LA, Strickland MJ, Sarnat SE, Sarnat JA,
\item Waller LA, Winquist A, Tolbert PE (2011). A Method for Detection of Residual Confounding in Time-series and Other Observational Studies. Epidemiology 22(1): 59-67.
\item Fleming TR, Harrington DP (2005). Counting Processes and Survival Analysis, 2nd ed. New York: Wiley.
\item Frangakis CE, Rubin DB (2002). Principal stratification in causal inference. Biometrics 58(1); 21-29.
\item Fulcher IR, Shpitser I, Marealle S, Tchetgen Tchetgen EJ (2020). Robust inference on population indirect causal effects: the generalized front door criterion. Journal of the Royal Statistical Society Series B: Statistical Methodology 82(1); 199-214.
\item Glymour MM, Spiegelman D (2016). Evaluating public health interventions: 5. Causal inference in public health research-Do sex, race, and biological factors cause health outcomes? American Journal of Public Health 107(1): 81-85.
\item Glymour MM, Tchetgen Tchetgen EJ, Robins JM (2012). Credible Mendelian randomization studies: approaches for evaluating the instrumental variable assumptions. American Journal of Epidemiology 175: 332-339.
\item Goodfellow I, Bengio Y, Courville A (2016). Deep Learning. Cambridge, MA: MIT Press.
\item Greenland S (1977). Response and follow-up bias in cohort studies. American Journal of Epidemiology 106(3):184-187.
\item Greenland S (1980). The effect of misclassification in the presence of covariates. American Journal of Epidemiology 112(4):564-569.
\item Greenland S (1987). Interpretation and choice of effect measures in epidemiologic analyses. American Journal of Epidemiology 125(5):761-768.
\item Greenland S (1991). On the logical justification of conditional tests for twoby-two contingency tables. The American Statistician 45(3):248-251.
\item Greenland S (1996a). Basic methods for sensitivity analysis of bias. International Journal of Epidemiology 25(6):1107-1116.
\item Greenland S (1996b). Absence of confounding does not correspond to collapsibility of the rate ratio or rate difference. Epidemiology 7(5):498-501.
\item Greenland S (2000). An introduction to instrumental variables for epidemiologists. International Journal of Epidemiology 29(4):722-729.
\item Greenland S (2003). Quantifying biases in causal models: classical confounding versus collider-stratification bias. Epidemiology 14:300-306.
\item Greenland S (2009a). Bayesian perspectives for epidemiologic research. III. Bias analysis via missing-data methods. International Journal of Epidemiology 38:1662-1673.
\item Greenland S (2009b). Relaxation penalties and priors for plausible modeling of nonidentified bias sources. Statistical Science 24:195-210.
\item Greenland S (2019). Some misleading criticisms of P-values and their resolution with S-values. The American Statistician 73(supplement 1): 106-114.
\item Greenland S, Brumback B (2002). An overview of relations among causal modelling methods. International Journal of Epidemiology 31:1030-1037.
\item Greenland S, Lash TL (2008). Bias analysis. In: Rothman KJ, Greenland S, Lash TL, eds. Modern Epidemiology, 3rd edition. Philadelphia, PA: Lippincott Williams \& Wilkins, pp. 345-380.
\item Greenland S, Lash TL, Rothman KJ (2008). Concepts of interaction. In: Rothman KJ, Greenland S, Lash TL, eds. Modern Epidemiology, 3rd edition. Philadelphia, PA: Lippincott Williams \& Wilkins, pp. 71-83.
\item Greenland S, Poole C (1988). Invariants and noninvariants in the concept of interdependent effects. Scandinavian Journal of Work, Environment \& Health 14(2):125–129.
\item Greenland S, Robins JM (1986). Identifiability, exchangeability, and epidemiological confounding. International Journal of Epidemiology 15:413-19.
\item Greenland S, Robins JM (1988). Conceptual problems in the definition and interpretation of attributable fractions. American Journal of Epidemiology 128:1185-1197.
\item Greenland S, Robins JM (2009). Identifiability, exchangeability, and confounding revisited. Epidemiologic Perspectives \& Innovations 6:4.
\item Greenland S, Pearl J, Robins JM (1999). Causal diagrams for epidemiologic research. Epidemiology 10:37-48.
\item Greenland S, Robins JM, Pearl J (1999). Confounding and collapsibility in causal inference. Statistical Science 14:29-46.
\item Greenland S, Rothman KJ (2008). Introduction to stratified analysis. In: Rothman KJ, Greenland S, Lash TL, eds. Modern Epidemiology, 3rd edition. Philadelphia, PA: Lippincott Williams \& Wilkins, pp. 258-282.
\item Grieve AP (2003). The number needed to treat: a useful clinical measure or a case of the Emperor's new clothes? Pharmaceutical Statistics 2:87-102.
\item Hajek J (1971). Comment on ``An essay on the logical foundations of survey sampling by D. Basu''. In: Godambe VP, Sprott DA, eds. Foundations of Statistical Inference. New York City, NY: Holt, Rinehart, and Winston; 1971 (p. 236).
\item Halloran ME, Struchiner CJ (1995). Causal inference in infectious diseases. Epidemiology 6: 142-151.
\item Hahn J, Todd P, van der Klaauw W (2001). Identification and estimation of treatment effects with a regression discontinuity design. Econometrica 69: 201-209
\item Hansen BB (2008). The prognostic analogue of the propensity score. Biometrika 95(2): 481-488.
\item Harrell FE (2015). Regression Modeling Strategies: With Applications to Linear Models, Logistic and Ordinal Regression, and Survival Analysis. New York: Springer.
\item Hartwig FP, Wang L, Davey Smith G, Davies NM (2023). Average causal effect estimation via instrumental variables: the no simultaneous heterogeneity assumption. Epidemiology 34(3):325-332.
\item Hastie TJ, Tibshirani RJ (1990). Generalized Additive Models. London: Chapman \& Hall.
\item Hastie TJ, Tibshirani RJ, Friedman J (2009). The Elements of Statistical Learning: Data Mining, Inference,and Prediction, 2nd ed. New York: Springer.
\item Heckman JJ, Vytlacil EJ (1999). Local instrumental variables and latent variable models for identifying and bounding treatment effects. Proceedings of the National Academy of Sciences USA 96:4730-4734.
\item Hernán MA (2010). The hazards of hazard ratios. Epidemiology 21(1):13-15.
\item Hernán MA (2016). Does water kill? A call for less casual causal inferences. Annals of Epidemiology 26: 674-680.
\item Hernán MA (2017). Selection bias without colliders. American Journal of Epidemiology 21(1):13-15.
\item Hernán MA, Brumback B, Robins JM (2001). Marginal structural models to estimate the joint causal effect of nonrandomized treatments. Journal of the American Statistical Association 96(454):440–448.
\item Hernán MA, Clayton D, Keiding N (2011). The Simpson's paradox unraveled. International Journal of Epidemiology 40:780-785.
\item Hernán MA, Cole SR, Margolick JB, Cohen MH, Robins JM (2005). Structural accelerated failure time models for survival analysis in studies with time-varying treatments. Pharmacoepidemiology and Drug Safety 14(7):477-491.
\item Hernán, Dahabreh IJ, Dickerman BA, Swanson SA (2025). The target trial framework for causal inference from observational data: Why and when is it helpful? Annals of Internal Medicine 178:402-407.
\item Hernán MA, Hernández-D´ıaz S, Werler MM, Mitchell AA (2002). Causal knowledge as a prerequisite for confounding evaluation: an application to birth defects epidemiology. American Journal of Epidemiology 155:176-184.
\item Hernán MA, Hernández-D´ıaz S, Robins JM (2004). A structural approach to selection bias. Epidemiology 15:615-625.
\item Hernán MA, Hernández-D´ıaz S (2012). Beyond the intention to treat in comparative effectiveness research. Clinical Trials; 9(1):48-55.
\item Hernán MA, Hsu J, Healy B (2019). A second chance to get causal inference right: a classification of data science tasks. Chance 32(1):42-49.
\item Hernán MA, Robins JM (2006a). Estimating causal effects from epidemiological data. Journal of Epidemiology and Community Health 60: 578-586.
\item Hernán MA, Robins JM (2006b). Instruments for causal inference: An epidemiologist's dream? Epidemiology 17(4): 360-372. Hernán MA, Alonso A, Logan R, Grodstein F, Michels KB, Willett WC,
\item Manson JE, Robins JM (2008). Observational studies analyzed like randomized experiments: an application to postmenopausal hormone therapy and coronary heart disease (with discussion). Epidemiology 19(6):766-779.
\item Hernán MA, Robins JM (2016). Using big data to emulate a target trial when a randomized trial is not available. American Journal of Epidemiology 183(8):758-764.
\item Hernán MA, Robins JM (2017). Per-protocol analyses of pragmatic trials. New England Journal of Medicine 377(14): 1391-1398.
\item Hernán MA, Sauer BC, Hernández-D´ıaz S, Platt R, Shrier I (2016). Specifying a target trial prevents immortal time bias and other self-inflicted injuries in observational analyses. Journal of Clinical Epidemiology 79: 70-75.
\item Hernán MA, Taubman SL (2008). Does obesity shorten life? The importance of well-defined interventions to answer causal questions. International Journal of Obesity 32: S8-S14.
\item Holland PW (1986). Statistics and causal inference (with discussion). Journal of the American Statistical Association 81:945-961.
\item Hoover DR, Muñoz A, He Y, Taylor JMG, Kingsley L, Chmiel JS, Saah A (1994). The effectiveness of interventions on incubation of AIDS as measured by secular increases within the population. Statistics in Medicine 13:2127-2139.
\item Horvitz DG, Thompson DJ (1952). A generalization of sampling without replacement from a finite universe. Journal of the American Statistical Association 47:663-685.
\item Huang Y, Valtorta M (2006). Pearl's calculus of intervention is complete. In: Proceedings of the 22nd Conference in Uncertainty in Artificial Intelligence. Cambridge, MA: AUAI Press, pp. 217-224.
\item Hosmer DW, Lemeshow S, May S (2008). Applied Survival Analysis: Regression Modelling of Time to Event Data. Hoboken, NJ: Wiley.
\item Hudgens MG, Halloran ME (2009). Towards causal inference with interference. Journal of the American Statistical Association 103:832-842. Hume D (1748). An Enquiry Concerning Human Understanding. Reprinted and edited 1993, Indianapolis/Cambridge: Hacket.
\item Imai K, Ratkovic M (2015). Robust estimation of inverse probability weights for marginal structural models. Journal of the American Statistical Association 110, 1013–1023.
\item Imbens GW (2004). Nonparametric estimation of average treatment effects under exogeneity: a review. The Review of Economics and Statistics 86(1): 4–29.
\item Imbens GW, Angrist JD (1994). Identification and estimation of local average treatment effects. Econometrica 62:467-475.
\item Imbens GW, Rubin DB (1997). Estimating outcome distributions for compliers in instrumental variables models. Review of Economic Studies 64:555-574.
\item Imbens G, Lemieux T (2008). Regression discontinuity designs: A guide to practice. Journal of Econometrics 142(2): 615–635.
\item Kalbfleisch and Prentice (2002). The Statistical Analysis of Failure Time Data. Hoboken, NJ: Wiley.
\item Kallus N, Santacatterina M (2018). Optimal balancing of time-dependent confounders for marginal structural models. arXiv preprint arXiv:1806.01083.
\item Katan MB (1986). Apolipoprotein E isoforms, serum cholesterol, and cancer. Lancet, 1:507-508.
\item Kleiner A, Talwalkar A, Sarkar P, Jordan MI (2014). A scalable bootstrap for massive data. Journal of the Royal Statistical Society B, 76 (Part 4):795-816.
\item Kosinski S, Stillwell D, Graepel T (2013). Private traits and attributes are predictable from digital records of human behavior. Proceedings of the National Academy of Sciences 110(15): 5802-5805. Kurth T, Walker AM, Glynn RJ, Chan KA, Gaziano JM, Berger K, Robins
\item JM (2006). Results of multivariable logistic regression, propensity matching, propensity adjustment, and propensity-based weighting under conditions of nonuniform effect. American Journal of Epidemiology 163(3): 262-270.
\item Korn EL, Baumrind S (1998). Clinician preferences and the estimation of causal treatment differences. Statistical Science 13:209-235
\item Laupacis A, Sackett DL, Roberts RS (1988). An assessment of clinically useful measures of the consequences of treatment. New England Journal of Medicine 318:1728-1733.
\item Lauritzen SL, Dawid AP, Larsen BN, Leimer H-G (1990). Independence properties of directed Markov fields. Networks 20:491-505.
\item Lash TL, Fox MP, Fink AK (2009). Applying Quantitative Bias Analysis to Epidemiologic Data. New York: Springer.
\item Lewis D (1973). Counterfactuals. Oxford: Blackwell.
\item Liang K-Y, Zeger SL (1986). Longitudinal data analysis using generalized linear models. Biometrika 73(1):13-22. Lin V, McGrath S, Zhang Z, Logan RW, Petito LC, Young JG, Hernán
\item MA (2019). gfoRmula: Parametric G-Formula. R package version 0.2.1. https://CRAN.R-project.org/package=gfoRmula.
\item Lin L, Mukherjee R, Robins JM (2020). On nearly assumption-free tests of nominal confidence interval coverage for causal parameters estimated by machine learning. Statistical Science 35(3): 518-539. Lipsitch M, Tchetgen Tchetgen E, Cohen T. Negative controls: a tool for detecting confounding and bias in observational studies. Epidemiology 21(3):383-8 (erratum in Epidemiology 2010;21(4):589).
\item Little RJ, D'Agostino R, Cohen ML, et al (2012). The prevention and treatment of missing data in clinical trials. New England Journal of Medicine 367(14): 1355-60.
\item Mackie JL (1965). Causes and conditions. American Philosohpical Quartely 2:245-264. Madenci AL, Kurgansky KE, Dickerman BA, Gerlovin H, Wanis KN, Smith AD, Trinqart L, Gagnon DR, Cho K, Gaziano JM, Casas JP, Robins JM, Hernán MA. Estimating the effect of bariatric surgery on cardiovascular events using observational data? Epidemiology 2024; 35(5):721-729.
\item Manski CF (1990). Nonparametric bounds on treatment effects. American Economic Review 80(2):319-323.
\item Marini MM, Olsen AR, Rubin DB (1980). Maximum-likelihood estimation in panel studies with missing data. Sociological Methodology 11:314-357.
\item Martens E, Pestman W, de Boer A, Belitser S, Klungel OH (2006). Instrumental variables: applications and limitations. Epidemiology 17(4):260-267.
\item McClellan M, McNeil BJ, Newhouse JP (2004). Does more intensive treatment of acute myocardial infarction in the elderly reduce mortality? Analysis using instrumental variables. JAMA 272(11):859-866.
\item McCullagh P, Nelder JA (1989). Generalized Linear Models, 2nd ed. London: Chapman \& Hall.
\item McCulloch CE, Searle SE, Neuhaus JM (2008). Generalized, Linear, and Mixed Models, 2nd ed. New York, NY: Wiley.
\item McGrath S, Young JG, Hernán MA (2022). Revisiting the g-null paradox. Epidemiology 33(1):114-120.
\item Meyer BD (1995). Natural and quasi-experiments in economics. Journal of Business \& Economic Statistics 13(2):151-161.
\item Miao W, Geng Z, Tchetgen Tchetgen EJ (2018). Identifying causal effects With proxy variables of an unmeasured confounder. Biometrika 105(4):987-993.
\item Miettinen OS (1972). Standardization of risk ratios. American Journal of Epidemiology 96:383-388.
\item Miettinen OS (1982). Causal and preventive interdependence: elementary principles. Scandinavian Journal of Work, Environment \& Health 8:159-168.
\item Miettinen OS, Cook EF (1981). Confounding: Essence and detection. American Journal of Epidemiology 1981; 114:593-603.
\item Molina J, Rotnitzky A, Sued M, Robins JM (2017) Multiple robustness in factorized likelihood models. Biometrika 104(3):561-581.
\item Neyman J (1923). On the Application of Probability Theory to Agricultural Experiments: Essay on Principles, Section 9. Translated in Statistical Science 1990; 5:465-480.
\item Ogburn EL, VanderWeele TJ (2012). On the nondifferential misclassification of a binary confounder. Epidemiology; 23(3):433-439.
\item Page J (2005). Doubly Robust Estimation: Structural Nested Cumulative Failure Time Models, Sc.D. dissertation, Departments of Epidemiology and Biostatistics, Harvard School of Public Health, Boston, MA. Palmer TM, Sterne JAC, Harbord RM, Lawlor DA, Sheehan NA, Meng S,
\item Granelli R, Davey Smith G, Didelez V (2011). Instrumental variable estimation of causal risk ratios and causal odds ratios in Mendelian randomization analyses. American Journal of Epidemiology 173(12): 1392-1403.
\item Pearl J (1988). Probabilistic Reasoning in Intelligent Systems. San Mateo, CA: Morgan Kaufmann.
\item Pearl J (1995). Causal diagrams for empirical research. Biometrika; 82:669-710.
\item Pearl J (2001). Direct and Indirect Effects. In: Breese J, Koller D, eds. Proceedings of the Seventeenth Conference on Uncertainty and Artificial Intelligence. San Francisco, CA: Morgan Kaufmann, pp. 411-420.
\item Pearl J (2009). Causality: Models, Reasoning, and Inference, 2nd ed. New York: Cambridge University Press.
\item Pearl J (2011). Understanding bias amplification. American Journal of Epidemiology 174(11):1223-1227.
\item Pearl J (2018). Does obesity shorten life? Or is it the soda? On nonmanipulable causes. Journal of Causal Inference 6(2); pp. 20182001.
\item Pearl J (2019). On the interpretation of do(x). Journal of Causal Inference 7(1); pp. 20192002.
\item Pearl J, Bareinboim (2014). External validity: From do-calculus to transportability Across Populations. Statistical Science 29(4): 579-595.
\item Pearl J, Robins JM (1995). Probabilistic evaluation of sequential plans from causal models with hidden variables. In: Proceedings of the 11th Conference on Uncertainty in Artificial Intelligence. Montreal, Canada, pp. 444-453. Pearson K, Lee A, Bramley-Moore L (1899). VI. Mathematical contributions to the Theory of Evolution.—VI. Genetic (Reproductive) selection: Inheritance of fertility in man, and of fecundity in thoroughbred horses. Philosophical Transactions of the Royal Society of London, Series A 192: 258-331.
\item Peters J, Janzing D, Schölkopf B (2017). Elements of Causal Inference: Foundations and Learning Algorithms. Cambridge, MA: MIT Press.
\item Petersen M, Schwab J, Gruber S, Blaser N, Schomaker M, van der Laan M (2014). Targeted maximum likelihood estimation for dynamic and static longitudinal marginal structural working models. Journal of Causal Inference 2(2):147-185.
\item Picciotto S, Hernán MA, Page J, Young JG, Robins JM (2012). Structural nested cumulative failure time models for estimating the effects of interventions. Journal of the American Statistical Association 107(499):886-900.
\item Richardson TS, Robins JM (2010). Analysis of the binary instrumental variable model. In: Dechter R, Geffner H, Halpern JY, eds. Heuristics, Probability and Causality: A Tribute to Judea Pearl. College Publications, UK.
\item Richardson TS, Robins JM (2014). ACE bounds; SEMs with equilibrium conditions. Statistical Science 29(3):363-366.
\item Richardson TS, Evans RJ, Robins JM (2010). Transparent parametrizations of models for potential outcomes. In: Bernardo JM, Bayarri MJ, Berger JO, Dawid AP, eds. Bayesian Statistics 9. Oxford University Press.
\item Richardson TS, Robins JM (2013). Single world intervention graphs (SWIGs): A unification of counterfactual and graphical approaches to causality. Working Paper 128, Center for Statistics and the Social Sciences, Seattle, WA.
\item Richardson TS, Robins JM, Wang L (2017). On modeling and estimation for the relative risk and risk difference. Journal of the American Statistical Association 112(519): 1121-1130.
\item Richardson WS, Wilson MC, Nishikawa J, Hayward RSA (1995). The wellbuilt clinical question: a key to evidence-based decisions. ACP Journal Club 123(3):A12-13.
\item Robins JM (1986). A new approach to causal Inference in mortality studies with sustained exposure periods -Application to control of the healthy worker survivor effect. Mathematical Modelling 7:1393-1512 (errata in Computers and Mathematics with Applications 1987;14:917-921).
\item Robins JM (1987). Addendum to ``A new approach to causal inference in mortality studies with sustained exposure periods -Application to control of the healthy worker survivor effect''. Computers and Mathematics with Applications 14(9-12):923-945 (errata in Computers and Mathematics with Applications 1987;18:477).
\item Robins JM (1988). Confidence intervals for causal parameters. Statistics in Medicine 7:773-785.
\item Robins JM (1989). The analysis of randomized and non-randomized AIDS treatment trials using a new approach to causal inference in longitudinal studies. In: Sechrest L, Freeman H, Mulley A, eds. Health Services Research Methodology: A Focus on AIDS. U.S. Public Health Service, National Center for Health Services Research, 113-159.
\item Robins JM (1993). Analytic methods for estimating HIV treatment and cofactor effects. In: Methodological Issues of AIDS Mental Health Research. Ostrow DG, Kessler R, eds. New York: Plenum Publishing, pp. 213-290.
\item Robins JM. (1994). Correcting for non-compliance in randomized trials using structural nested mean models. Communications in Statistics 23:2379-2412.
\item Robins JM (1997a). Causal Inference from Complex Longitudinal Data. Latent Variable Modeling and Applications to Causality. Berkane M, ed. New York, NY: Springer Verlag, pp. 69-117.
\item Robins JM (1997b). Structural nested failure time models. In: Survival Analysis, Andersen PK, Keiding N, Section Editors. Encyclopedia of Biostatistics, Armitage P, Colton T (eds). Chichester, UK: John Wiley \& Sons, 4372-4389.
\item Robins JM (1998a). Marginal structural models. 1997 Proceedings of the Section on Bayesian Statistical Science. Alexandria, Virginia: American Statistical Association, 1-10.
\item Robins JM (1998b). Correction for non-compliance in equivalence trials. Statistics in Medicine 17: 269–302.
\item Robins JM (1999). Marginal structural models versus structural nested models as tools for causal inference. In: Halloran E, Berry D. Statistical Models in Epidemiology: The Environment and Clinical Trials. New York, Springer-Verlag: 95-134.
\item Robins JM (2000). Robust estimation in sequentially ignorable missing data and causal inference models. 1999 Proceedings of the Section on Bayesian Statistical Science. Alexandria, Virginia: American Statistical Association, pp. 6-10.
\item Robins JM (2001). Data, design, and background knowledge in etiologic inference. Epidemiology 11(3):313-320.
\item Robins JM, Greenland S (1986). The role of model selection in causal inference from nonexperimental data. American Journal of Epidemiology 123(3): 392-402.
\item Robins JM, Greenland S (1989). Estimability and estimation of excess and etiologic fraction. Statistics in Medicine 8:845-859.
\item Robins JM, Greenland S (1992). Identifiability and exchangeability for direct and indirect effects. Epidemiology 3(2):143-155.
\item Robins JM, Greenland S (2000). Comment on ``Causal inference without counterfactuals.'' Journal of the American Statistical Association 95:477-82.
\item Robins JM, Hernán MA, Rotnitzky A (2007). Effect modification by timevarying covariates. American Journal of Epidemiology 166:994-1002.
\item Robins JM, Hernán MA, Siebert U (2004). Effects of multiple interventions. In: Comparative Quantification of Health Risks: Global and Regional Burden of Disease Attributable to Selected Major Risk Factors Vol II. Ezzati M, Lopez AD, Rodgers A, Murray CJL, eds. Geneva: World Health Organization, 2004.
\item Robins JM, Morgenstern H (1987). The foundations of confounding in epidemiology. Computers \& Mathematics with Applications 14(9-12): 869-916.
\item Robins JM, Orellana L, Rotnitzky A (2008). Estimation and extrapolation of optimal treatment and testing strategies. Statistics in Medicine 27(23):4678-721.
\item Robins JM, Richardson TS (2010). Alternative graphical causal models and the identification of direct effects. In: Causality and Psychopathology: Finding the Determinants of Disorders and Their Cures. P. Shrout, ed. New York, NY: Oxford University Press.
\item Robins JM, Richardson TS, Shpitser I (2022). An interventionist approach to mediation analysis. In: Probabilistic and Causal Inference: The Works of Judea Pearl. Geffner H, Dechter R, Halpern JY, eds. New York, NY: Association for Computing Machinery.
\item Robins JM, Ritov Y (1997). Toward a curse of dimensionality appropriate (CODA) asymptotic theory for semiparametric models. Statistics in Medicine 17:285-319.
\item Robins JM, Rotnitzky A (1992). Recovery of information and adjustment for dependent censoring using surrogate markers. In: AIDS Epidemiology – Methodological Issues. Jewell N, Dietz K, Farewell V, eds. Boston, MA: Birkhäuser, pp. 297-331.
\item Robins JM, Rotnitzky A, Scharfstein D (1999). Sensitivity analysis for selection bias and unmeasured confounding in missing data and causal inference models. In: Halloran ME, Berry D, eds. Statistical Models in Epidemiology: The Environment and Clinical Trials. New York, NY: Springer-Verlag, pp. 1-92.
\item Robins JM, Rotnitzky A, Vansteelandt S (2007). Discussion on ``Principal stratification designs to estimate input data missing due to death.'' Biometrics 63(3):650-654.
\item Robins JM, Rotnitzky A, Zhao LP (1994) Estimation of regression coefficients when some regressors are not always observed. Journal of the American Statistical Association 89: 846-866
\item Robins JM, Scheines R, Spirtes P, Wasserman L (2003). Uniform consistency in causal inference. Biometrika 90(3):491-515.
\item Robins JM, Wasserman L (1997). Estimation of effects of sequential treatments by reparameterizing directed acyclic graphs. In: Geiger D, Shenoy P, eds. Proceedings of the Thirteenth Conference on Uncertainty in Artificial Intelligence. San Francisco, CA: Morgan Kaufmann, pp. 409-420.
\item Robins JM, Weissman M (2016). Counterfactual causation and streetlamps. What is to be done? International Journal of Epidemiology 45(6):1830-1835.
\item Rosenbaum PR (1984). The consequences of adjustment for a concomitant variable that has been affected by the treatment. Journal of the Royal Statistical Society Series A 147:656-666.
\item Rosenbaum PR (1987). Model-based direct adjustment. Journal of the American Statistical Association 82:387-394.
\item Rosenbaum PR (2007). Interference between units in randomized experiments. Journal of the American Statistical Association 102:191-200.
\item Rosenbaum PR (2002). Observational Studies, 2nd edition. New York, NY: Springer-Verlag.
\item Rosenbaum PR (2005). Sensitivity analysis in observational studies. In: Encyclopedia of Statistics in Behavioral Sciences, Everitt BS, Howell DC T (eds). Chichester, UK: John Wiley \& Sons, 4:1809-1814.
\item Rosenbaum PR, Rubin DB (1983). The central role of the propensity score in observational studies for causal effects. Biometrika 70:41-55.
\item Rothman KJ (1976). Causes. American Journal of Epidemiology 104:587-592.
\item Rothman KJ, Greenland S, Walker AM (1980). Concepts of interaction. American Journal of Epidemiology 112:467-470.
\item Rothman KJ, Greenland S, Poole C, Lash TL (2008). Causation and Causal Inference. In: Rothman KJ, Greenland S, Lash TL, eds. Modern Epidemiology, 3rd ed. Philadelphia, PA: Lippincott Williams and Wilkins.
\item Rubin DB (1974). Estimating causal effects of treatments in randomized and nonrandomized studies. Journal of Educational Psychology 56:688-701.
\item Rubin DB (1976). Inference and missing data (with discussion). Biometrika 63:581-592.
\item Rubin DB (1978). Bayesian inference for causal effects: The role of randomization. Annals of Statistics 6:34-58.
\item Rubin DB (1980). Discussion of ``Randomized analysis of experimental data: the Fisher randomization test'' by Basu D. Journal of the American Statistical Association 75:591-593.
\item Rubin DB (2004). Direct and indirect effects via potential outcomes. Scandinavian Journal of Statistics 31(2):161-170.
\item Rudolph KE, van der Laan MJ (2017). Robust estimation of encouragementdesign intervention effects transported across sites. Journal of the Royal Statistical Society Series B 79(5):1509-1525.
\item Samuels ML (1981). Matching and design efficiency in epidemiological studies. Biometrika 68:577-588.
\item Santosa F, Symes WW (1986). Linear inversion of band-limited reflection seismograms. SIAM Journal on Scientific and Statistical Computing 7(4): 1307-1330.
\item Saracci R (1980). Interaction and synergism. American Journal of Epidemiology 112:465-466.
\item Sato T, Matsuyama Y (2003). Marginal structural models as a tool for standardization. Epidemiology 14:680-686.
\item Scharfstein DO, Rotnitzky A, Robins JM (1999). Adjusting for nonignorable drop-out using semiparametric nonresponse models. Journal of the American Statistical Association 94(448): 1096-1120.
\item Schwartz S, Gatto NM, Campbell UB (2016). Causal identification: a charge of epidemiology in danger of marginalization. Annals of Epidemiology 26(10):669-673.
\item Shahn Z, Dukes O, Richardson D, Tchetgen Tchetgen EJ, Robins JM (2022). Structural nested mean models under parallel trends assumptions. arXiv preprint arXiv:2204.10291.
\item Shpitser I, Pearl J (2006). Identification of joint interventional distributions in recursive semi-Markovian causal models. In: Proceedings of the 21st National Conference on Artificial Intelligence. Menlo Park, CA: AAAI Press, pp. 1219-1226.
\item Shpitser I, Richardson TS, Robins JM (2022). Multivariate counterfactual systems and causal graphical models. In: Probabilistic and Causal Inference: The Works of Judea Pearl. Geffner H, Dechter R, Halpern JY, eds. New York, NY: Association for Computing Machinery.
\item Shpitser I, Tchetgen Tchetgen EJ, Andrews R (2021). Modeling interference via symmetric treatment decomposition. arXiv preprint arXiv:1709.01050
\item Simpson EH (1951). The interpretation of interaction in contingency tables. Journal of the Royal Statistical Society, Series B 13:238-241.
\item Smith GCS, Pell JP (2003). Parachute use to prevent death and major trauma related to gravitational challenge: systematic review of randomised controlled trials. British Medical Journal 327:1459-1461.
\item Sobel ME (2006). What do randomized studies of housing mobility demonstrate? Journal of the American Statistical Association 101: 1398-1407.
\item Sofer T, Richardson DB, Colicino E, Schwartz J, Tchetgen Tchetgen EJ (2016). On negative outcome control of unobserved confounding as a generalization of difference-in-differences. Statistical Science 31(3): 348-361.
\item Sommer A, Zeger SL (1991). On estimating efficacy from clinical trials. Statistics in Epidemiology 10:45-52.
\item Spirtes P, Glymour C, Scheines R (2000). Causation, Prediction and Search, 2nd ed. Cambridge, MA: MIT Press.
\item Stalnaker RC (1968). A theory of conditionals. In Rescher N, ed. Studies in Logical Theory. Oxford: Blackwell. Reprinted in Jackson F, ed. Conditionals. Oxford: Oxford University Press, 1991.
\item Stensrud MJ, Young JG, Didelez V, Robins JM, Hernán MA (2020). Separable effects for causal inference in the presence of competing risks. Journal of the American Statistical Association 117(537):175-183. Stensrud MJ, Hernán MA, Tchetgen Tchetgen EJ, Robins JM, Didelez V,
\item Young JG (2021). A generalized theory of separable effects in competing event settings. Lifetime Data Analysis 27(4):588-631.
\item Stensrud MJ, Robins JM, Sarvet A, Tchetgen Tchetgen EJ, Young JG (2023). Conditional separable effects. Journal of the American Statistical Association (in press).
\item Stuart EA (2010). Matching methods for causal inference. Statistical Science 25, 1-21.
\item Swanson SA, Hernán MA (2013). How to report instrumental variables analyses (suggestions welcome). Epidemiology 24(3): 370-374.
\item Swanson SA, Hernán MA (2014). Think globally, act globally: An epidemiologist's perspective on instrumental variable estimation. Statistical Science 29(3): 371-374.
\item Swanson SA, Hernán MA, Miller M, Robins JM, Richardson T (2018). Partial identification of the average treatment effect using instrumental variables: Review of methods for binary instruments, treatments, and outcomes. Journal of the American Statistical Association 113(522):933-947. Swanson SA, Holme Ø, Løberg M, Kalager M, Bretthauer M, Hoff G, Aas E, Hernán MA (2015a). Bounding the per-protocol effect in randomized trials: an application to colorectal cancer screening. Trials 16:541.
\item Swanson SA, Miller M, Robins JM, Hernán MA (2015b). Definition and evaluation of the monotonicity condition for preference-based instruments. Epidemiology 26:414-420.
\item Swanson SA, Miller M, Robins JM, Hernán MA (2015c). Selecting on treatment: a pervasive form of bias in instrumental variable analysis. American Journal of Epidemiology 181(3):191-197.
\item Tchetgen Tchetgen EJ (2009). A commentary on G. Molenberghs's review of missing data methods. Drug Information Journal 43(4):433–435.
\item Tchetgen Tchetgen EJ, Rotnitzky A (2011). Double-robust estimation of an exposure-outcome odds ratio adjusting for confounding in cohort and case-control studies. Statistics in Medicine 30(4):335-47. Tchetgen Tchetgen EJ, Stefan W, Vansteelandt S, Martinussen T, Glymour M (2015). Instrumental variable estimation in a survival context. Epidemiology 26(3): 402-410.
\item Thistlewaite D, Campbell D (1960). Regression-discontinuity analysis: An alternative to the ex-post facto experiment. Journal of Educational Psychology 51:309-317.
\item Thompson WA (1977). On the treatment of grouped observations in life studies. Biometrics 33:463-470.
\item Tian J, Pearl J (2002). A general identification condition for causal effects. In: Proceedings of the 18th National Conference on Artificial Intelligence. Menlo Park, CA: AAAI Press/MIT Press, pp. 567–573.
\item Tibshirani R (1996). Regression shrinkage and selection via the lasso. Journal of the Royal Statistical Society Series B 58(1): 267-288. van der Laan MJ, Petersen ML (2007). Causal effect models for realistic individualized treatment and intention to treat rules. International Journal Biostatistics 3(1): Article 3. van der Laan MJ, Rubin D (2006). ) Targeted maximum likelihood learning. International Journal of Biostatistics 2(1): Article 11. van der Laan MJ, Gruber S (2012). Targeted minimum loss based estimation of causal effects of multiple time point interventions. International Journal of Biostatistics 8(1): Article 9.
\item VanderWeele TJ (2009a). Concerning the consistency assumption in causal inference. Epidemiology 20:880-883.
\item VanderWeele TJ (2009b). On the distinction between interaction and effect modification. Epidemiology 20:863-871.
\item VanderWeele TJ (2010a). Empirical tests for compositional epistasis. Nature Reviews Genetics 11:166.
\item VanderWeele TJ (2010b). Sufficient cause interactions for categorical and ordinal exposures with three levels. Biometrika 97:647-659.
\item VanderWeele TJ (2015). Explanation in Causal Inference. Methods for Mediation and Interaction. New York, NY: Oxford University Press.
\item VanderWeele TJ, Hernán MA (2006). From counterfactuals to sufficient component causes and vice versa. European Journal of Epidemiology 21:855-858.
\item VanderWeele TJ, Hernán MA (2009). Results on differential and dependent measurement error of the exposure and the outcome using signed directed acyclic graphs. American Journal of Epidemiology 175(12):1303-1310.
\item VanderWeele TJ, Hernán MA, Robins JM (2008). Causal directed acyclic graphs and the direction of unmeasured confounding bias. Epidemiology 19:720-728.
\item VanderWeele TJ, Arah OA (2011). Bias formulas for sensitivity analysis of unmeasured confounding for general outcomes, treatments, and confounders. Epidemiology 22:42-52.
\item VanderWeele TJ, Robins JM (2007a). The identification of synergism in the sufficient-component-cause framework. Epidemiology 18:329-339.
\item VanderWeele TJ, Robins JM (2007b). Four types of effect modification. A classification based on directed acyclic graphs. Epidemiology 18; 561-568.
\item VanderWeele TJ, Robins JM (2007c). Directed acyclic graphs, sufficient causes, and the properties of conditioning on a common effect. American Journal of Epidemiology 166; 1096-1104.
\item VanderWeele TJ, Robins JM (2008). Empirical and counterfactual conditions for sufficient cause interactions. Biometrika 95:49-61.
\item VanderWeele TJ, Robins JM (2012). Stochastic counterfactuals and stochastic sufficient causes. Statistica Sinica 22:379-392.
\item VanderWeele TJ, Shpitser I (2013). On the definition of a confounder. Annals of Statistics 41(1): 196-220.
\item VanderWeele TJ, Tchetgen Tchetgen EJ, Cornelis M, Kraft P (2014). Methodological challenges in Mendelian randomization. Epidemiology 25(3):427-435.
\item Vansteelandt S, Goetghebeur E (2003). Causal inference with generalized structural mean models. Journal of the Royal Statistical Society - Series B 65;817-835.
\item Wald A (1940). The fitting of straight lines if both variables are subject to error. Annals of Mathematical Statistics 11:284-300.
\item Walker AM (1991). Observation and Inference: An Introduction to the Methods of Epidemiology. Newton Lower Falls, MA: Epidemiology Resources, Inc.
\item Wang L, Meng X, Richardson TS, Robins JM (2022). Coherent modeling of longitudinal causal effects on binary outcomes. Biometrics (in press).
\item Wang L, Tchetgen Tchetgen E (2018). Bounded, efficient and multiply robust estimation of average treatment effects using instrumental variables. Journal of the Royal Statistical Society Series B 80 (Part 3):531-550.
\item Wang Y, Zubizarreta JR (2020). Minimal dispersion approximately balancing weights: Asymptotic properties and practical considerations. Biometrika 107, 93–105.
\item Wasserman L (2004). All of Statistics: A Concise Course in Statistical Inference. New York, NY: Springer.
\item Weinberg CR, Umbach DM, Greenland S (1994). When will nondifferential misclassification of an exposure preserve the direction of a trend? American Journal of Epidemiology 140:565-571.
\item Wen L, Young JG, Robins JM, Hernán MA (2021). Parametric g-formula implementations for causal survival analyses. Biometrics 77(2):740-753.
\item Wen L, Hernán MA, Robins JM (2022). Multiply robust estimators of causal effects for survival outcomes. Scandinavian Journal of Statistics 49:1304-1328.
\item Wen L, Sarvet AL, Stensrud MJ (2023). Causal effects of intervening variables in settings with unmeasured confounding. arXiv preprint arXiv:2305.00349.
\item Westreich D, Edwards JK, Lesko CR, Stuart E, Cole SR (2017). Transportability of trial results using inverse odds of sampling weights. American Journal of Epidemiology 186(8):1010-1014.
\item Wold H (1954). Causality and econometrics Econometrica 22(2):162-177.
\item Wooldridge JM (2010). Econometric Analysis of Cross section and Panel Data. Cambridge, MA: MIT Press, 2nd edition.
\item Young JG, Cain LE, Robins JM, O'Reilly E, Hernán MA (2011). Comparative effectiveness of dynamic treatment regimes: an application of the parametric g-formula. Statistics in Biosciences 3:119-143.
\item Young JG, Hernán MA, Robins JM (2014). Identification, estimation and approximation of risk under interventions that depend on the natural value of treatment using observational data. Epidemiologic Methods 3(1):1-19.
\item Young JG, Stensrud MJ, Tchetgen Tchetgen EJ, Hernán MA (2019). A causal framework for classical statistical estimands in failure time settings with competing events. Statistics in Medicine 39(8):1199-1236.
\item Yule GU (1903). Notes on the theory of association of attributes in statistics. Biometrika 2:121-134.
\end{list}
\endgroup
